\documentclass[a4paper,10pt]{article}

\usepackage[margin=2cm]{geometry}
\usepackage{graphicx}
\usepackage{amsmath}
\usepackage{array}
\usepackage{hyperref}
\usepackage[all]{hypcap}
\usepackage{listings}
\lstdefinestyle{TerminalStyle}{
  language=bash,
  basicstyle=\small\sffamily,
  numbers=left,
  numberstyle=\tiny,
  numbersep=3pt,
  frame=tb,
  columns=fullflexible,
  linewidth=0.9\linewidth,
  xleftmargin=0.1\linewidth
}
\lstdefinestyle{HtmlStyle}{
  language=html,
  basicstyle=\small\sffamily,
  numbers=left,
  numberstyle=\tiny,
  numbersep=3pt,
  frame=tb,
  columns=fullflexible,
  linewidth=0.9\linewidth,
  xleftmargin=0.1\linewidth
}
\lstdefinestyle{OutputStyle}{
  language=html,
  basicstyle=\small\sffamily,
  frame=tb,
  columns=fullflexible,
  linewidth=0.9\linewidth,
  xleftmargin=0.1\linewidth
}

\setlength{\parindent}{0pt}
\setlength{\parskip}{1ex plus 0.5ex minus 0.2ex}
\title{\includegraphics[width=12cm]{Eeufeeslogo.jpg} \\
       Department of Computer Science \\
       University of Pretoria \\
       \vspace{0.5cm}
       Software Engineering\\
       COS301 MiniProject \\
       \vspace{0.5cm}
       \begin{large} \textbf{Team Maroon}\\ NavUP\end{large}}

\date{} 
\author{Bondjobo, Jocelyn (J) 	13232852 		\\
		Mweshi, George (G)		16394713		\\
		Letsoalo, Joseph (J)	15043844		\\
		Setoaba, Phuti (P)		13032616		\\
		Trivella, Camron (C)	14070970		\\
		Coetzer, Albert (A)		15244882		\\
}

\begin{document}
\maketitle
\thispagestyle{empty}
\clearpage

\newpage
\pagenumbering{roman}
\thispagestyle{empty}
\tableofcontents
\clearpage

\newpage

\pagenumbering{arabic}

\section{Introduction}

	\subsection{Purpose} 	\subsection{Scope} 
	\subsection{Overview} 

\section{Overall Description}

	\subsection{Product Perspective}
	
		\subsubsection{System Interfaces}
			Each component and system/subsystem in the overall program will need its own interface between its own\\				software and the bigger program. For instance, the crowd sourcing routines must be able to record and report a\\			density map to the main branch of the program that must then be able to be accessed and utilized by an analyzer\\			which must then return it to the main branch to be retrieved as needed for navigation by the user.\\
		\subsubsection{User Interfaces}
			The user interface must consist of an easily understandable and navigable GUI, simple enough for even the\\				most inexperienced to use. The different options and account information must be visible and responsive and the\\			navigation component must be clear and accurate. The user should intuitively know how to get from any point to\\			any other point in the program.\\
		\subsubsection{Hardware Interfaces}
			The hardware interface must allow the program to access and utilize systems on the device on which the program\\			is running, specifically the GPS system if the program uses it, the wireless network adapter to check for Wi-Fi\\			routers and signal strength, the screen to display information and the device used to make selections. The\\				hardware interface will be the backbone of the program, allowing directions and instructions to be communicated\\			to the user.\\
		\subsubsection{Software Interfaces}
			The routines by which the program will be able to call, run and interpret other software on the device being\\				used to run the program. The routines will, for example, utilize the software responsible for managing the\\ 				connection between the device and the Wi-Fi routers. This will co-exist with the hardware interface to make use\\			of the hardware on the device for the program.\\
		\subsubsection{Communications Interfaces}
			The communications interface will most likely be used to connect devices being used to run the program to\\				other devices being used to run the program. One of the uses for this will be in crowd sourcing density data of\\			the users to help effectively and intelligently create the most efficient instructions for a specific user or to\\			record and utilize the most common trends in user movements to predict behavior and increase efficiency.\\
		\subsubsection{Memory}
			Information such as maps, router positions and user identities will most likely need to be stored on a server\\ 			by the program administrators. Information gathered about the users, such as common routes, time spent using\\ 				the program, user information and any other data collected by the program should also be stored in the same\\ 				way.However, an alternative would be to store the information on the user’s device and report it to the server\\ 			when needed.\\   
		\subsubsection{Operations}
			The primary operation of the program will be directing students around the University of Pretoria’s main\\				campus, secondary operations may include crowd sourcing utilities and a rewards program. It should be possible\\			to store and retrieve information about a user’s movements in order to back trace if necessary.\\
		\subsubsection{Site Adaptation Requirements}
			The program should be available in a variety of language options.\\
		
	\subsection{Product Functions}
		1. Basic Functionality\\
			NavUp has basic funtionality that allows the user's current location to be determined indoors and outdoors. The user can further search for a location, save a location, navigate to a location and see the route in 3D. Zoom in or out capability is also available. Location markers  are there to help the user differentiate their starting point, destination as well as other locations. The application can also generate an time estimate that the user may take to get to their destination.\\
		\\
		2. Surveillance\\
			Through surveillance, NavUP is able to provide and visualise pedestrian traffic information on campus and further provide an alternative route if one is available to avoid pedestrian traffic. Another capability that can be achieved is to record the user's number of steps which can be used by the user to track their movement for fitness or leisure or by a third party to run competitions.\\
		\\
		3. Extra features\\
			The appliaction can give recommendations of locations on campus in accordance to the user’s interests and preferences. And Rate and review places on campus\\
			
	\subsection{User Characteristics}
		The application NavUp can expect to have three types of users; the registered student or staff member, the guest and the administrator.\\
		\\
			1. The registered student or staff member and guest user are expected to be able to provide a destination location in terms of name of the building or faculty and room.\\
			\\
			2. The registered student or staff member will be able to save a track record of movements, favourite locations.\\
			\\
			3. The guest will have limited functionality; hence will only be able to navigate from place to place.\\
			\\
			4. The administrator must be able to set goals and provide information regarding competitions to be facilitated through NavUp. This entails the award giving process either for attendance of an event or achieving the set number of steps.\\

	\subsection{Constraints}
	\subsection{Assumptions and Dependencies}

\section{Specific Acquirements}

	\subsection{External Interface Requirements}
	\subsection{Functional Requirements}
	\subsection{Performance Requirements}
	\subsection{Design Constraints		\\
		1. User Movement		\\
			Because user location changes, to account for the location of the user in a building with one/ weak  \\
			Wi-Fi connection will lead to failure in updating to user location.		\\
		2. Failure connection		\\
			Since the app uses Wi-Fi to find location and navigate if the user fails to connect to Tuks Wi-Fi \\
			at that moment the app will not work regardless of whether the user uses cellular network data.		\\
		3. Location in buildings 	\\
			Finding the current location of the user and displaying it on the app inside the buildings will \\
			be restricted since we don’t have the lay-out/map of all buildings.	\\
			
	}
	\subsection{Software System Attributes}
	\subsection{Other Requirements}
	\subsubsection{Technology and Android clients requirements}
\begin{enumerate}
 \item \textbf{Programming Languages}
	\begin{itemize}
		\item Java
		\item eXtensible Markup Language (XML)
	\end{itemize}
\item \textbf{Frameworks}
	\begin{itemize}
		\item Android Studio
	\end{itemize}
\item \textbf{Libraries}
	\begin{itemize}
		\item Android Butterknife
	\end{itemize}
\item \textbf{Database System} \\ \\ 
Couchbase Mobile which consists of:
	\begin{itemize}
 		\item Couchbase Lite
		\item Couchbase Sync Gateway
		\item Couchbase Server
	\end{itemize}
\item \textbf{Operating System}
	\begin{itemize}
 		\item Android 4.0. Devices and upwards
	\end{itemize}
\item \textbf{Dependency Management and Build Tools}
	\begin{itemize}
		\item Gradle
	\end{itemize}

\end{enumerate}
	

\clearpage
\section{Open Issues}
\subsection {GitHub Repository}
Team Maroon Repository: \url{https://github.com/josephbondjobo/COS301-Maroon-Team/tree/develop/Task%201}

This repository contains:
\begin{enumerate}
\item All work done by team members.
\item \href{https://github.com/josephbondjobo/COS301-Maroon-Team/blob/develop/Intro/Contributors.md}{CONTRIBUTORS.md} file outlining which members where involved in which phases of the project.
\end{enumerate}



\newpage
\clearpage
\addcontentsline{toc}{section}{References}

\end{document}
